\documentclass{article}
\usepackage{graphicx} % Required for inserting images
\usepackage{amsmath}
\usepackage{amsthm}
\usepackage{amssymb}
\usepackage{mathrsfs}
\newtheorem{theorem}{Theorem}
\newtheorem{lemma}{Lemma}
\newtheorem{corollary}{Corollary}
\theoremstyle{remark}\newtheorem{remark}{Remark}
\newcommand{\T}[0]{\mathcal{T}}
\newcommand{\B}[0]{\mathcal{B}}
\newcommand{\C}[0]{\mathcal{C}}
\newcommand{\R}[0]{\mathbb{R}}
\newcommand{\Q}[0]{\mathbb{Q}}
\newcommand{\Z}[0]{\mathbb{Z}}
\newcommand{\N}[0]{\mathbb{N}}
\newcommand{\e}[0]{\epsilon}

\begin{document}

\noindent 2.28 For a formula $\phi$ built up using the connectives $\neg, \wedge, \vee$, let $\phi^*$ be constructed by replacing each propositional variable in $\phi$ by its negation. So if $\phi$ is the formula $((q \vee p) \wedge \neg p)$, $\phi^*$ is $((\neg q \vee \neg p) \wedge \neg\neg p)$.

\begin{enumerate}
    \item[(a)] For any truth assignment $v$, let $v^*$ be the truth assignment which gives each propositional variable the opposite value to that given by $v$, i.e.
    \[
        v^*(p) =
        \begin{cases}
            T, & \text{if } v(p) = F, \\
            F, & \text{if } v(p) = T,
        \end{cases}
    \]
    for all propositional variables $p$. Show that $v(\phi) = v^*(\phi^*)$. [\emph{Hint}: This is really a statement about all formulas $\phi$ of a certain sort, so what is the likely method of proof?]

    \item[(b)]
    \begin{enumerate}
        \item[(i)] Use the result of part (a) to show that $\phi$ is a tautology if and only if $\phi^*$ is a tautology.
        \item[(ii)] Is it true that $\phi$ is a contradiction if and only if $\phi^*$ is a contradiction? Explain your answer.
    \end{enumerate}
\end{enumerate}

\begin{proof}[proof of (a)]
    We proceed by induction on the length of a formula

    \textit{Base Case}: Suppose we have a formula, $\phi$, of length zero. So then $\phi$ is just a propositional variable, $p$. WLOG, let $v(p) = T$. Then $\phi^* = \neg p$. So then $v^* (\phi^*) = v^*(\neg p) = T$.

    \textit{Induction Hypothesis}: For all formulas, $\phi$, of length $\leq n$, $v(\phi) = v^*(\phi^*)$.

    \textit{Inductive Step}: Let $\phi$ be a formula of length $n + 1$. Since $n + 1 \geq 0$, we know that $\phi$ contains connective, so it isn't a propositional variable, so by the definition of formula it is one of these forms: $\neg \theta, (\theta \land \psi), (\theta \lor \psi)$.

    \textit{Case $\phi = \neg \theta$}: 
    
    Then $\theta$ has length $\leq n$, so by the inductive hypothesis $v(\theta) = v^*(\theta^*)$.

    By the definition of $*$, $\phi^* = \neg (\theta^*)$, so we want to show $v(\neg\theta) = v^* (\neg(\theta^*))$.
    
    Since the truth table for $\neg$ is symmetric under interchanging $T$ and $F$, WLOG suppose that $v(\theta) = v^* (\theta^*) = F$.
    
    So then $v(\phi) = v(\neg \theta) = T$.

    Since $v^* (\theta^*) = F$, $v^*(\neg (\theta^*)) = T = v^*(\phi^*)$.

    And so, $v(\neg\theta) = v^* (\neg(\theta^*))$, i.e., $v(\phi) = v^*(\phi^*)$.

    %-----------------------------------------------------

    % \textit{Case $\phi = (\theta \land \psi)$}:
    
    % Then $\theta$ and $\psi$ have length $\leq n$, so by the inductive hypothesis $v(\theta) = v^*(\theta^*)$ and $v(\psi) = v^*(\psi^*)$.

    % By the definition of $*$, $\phi^* = (\theta^* \land \psi^*)$, so we want to show $v((\theta \land \psi)) = v^*((\theta^* \land \psi^*))$.


    %-----------------------------------------------------
    
    \textit{Case $\phi = (\theta \lor \psi)$}: 

    Then $\theta$ and $\psi$ have length $\leq n$, so the the inductive hypothesis $v(\theta) = v^*(\theta^*)$ and $v(\psi) = v^*(\psi^*)$.

    By the definition of $*$, $\phi^* = (\theta^* \lor \psi^*)$, so we want to show $v((\theta \lor \psi)) = v^*((\theta^* \lor \psi^*))$.

    Since the truth table for $\lor$ is symmetric under interchanging $T$ and $F$, WLOG suppose that $v(\theta) = v^* (\theta^*) = T$ and $v(\psi) = v^* (\psi^*) = F$.

    So then $v(\phi) = v((\theta \lor \psi)) = T$.

    But (this is where i get confused)
    $v^*(\phi^*) = v^*((\theta^* \lor \psi^*)) = F$
    
    
\end{proof}

\end{document}
